\section{Conclusion}

Official exit is not institutional exit. A local government financing vehicle
may leave a regulatory list or adopt no-government-financing language while the
same public-finance function remains with the issuer, moves to another
state-owned entity, enters formal fiscal channels, or disappears with the legal
entity. This paper makes that institutional fate measurable by pairing a
documented formal event with evidence about the post-event location of the
function and assigning one of four mutually exclusive labels.

The measurement framework is the paper's main contribution. The current
working-reference file contains 94 Codex source-packet-reviewed cases awaiting
independent human confirmation, each linked to an event,
function evidence, an alternative label, and quality scores. All evidence
identifiers resolve to the source inventory, 93 cases have complete local
copies, and all 94 have recovery URLs. These features make each decision
auditable. They do not substitute for independent human validation: no human
reviewer has yet frozen confirmation decisions, and the current cases were not drawn from a
national probability frame.

The LLM screen serves the same measurement goal at a larger scale. It finds
one-sided nominal-exit candidates in repeated bond disclosures, but it does not
supply final four-category outcomes. The 61 issuers shared by the screen and
reference set agree exactly, yet that selected overlap cannot identify
population precision, recall, or multiclass error. Design-based supervised
learning can support validation-adjusted inference only after the project
verifies the issuer frame, draws cases with known probabilities, and obtains
independent human labels.

The historical application is therefore evidence about how the measure can be
used, not the paper's settled causal result. Among 84 matched reference cases,
historical elite density is positively associated with movement away from
nominal exit in simple models. The adjusted association is fragile under
complete-case restriction and case deletion, and the proxy may capture
education or social capital as well as administrative development.
Documentary patterns are consistent with fiscal absorption, bureaucratic
coordination, and project or state-asset governance, but the current design
does not identify those processes as mediators.

The resulting contribution is narrower and more durable than an official exit
count. It provides a source-auditable way to distinguish formal compliance from
functional fiscal change and shows why that distinction matters for the study
of central mandates, local adaptation, and quasi-fiscal organizations. A
submission-ready causal application will require a verified probability frame,
an independent second-coder round, broader representation of rare institutional
fates, and stable contemporary controls. Those additions extend the measure;
they do not change the institutional distinction on which the paper is built.
